\begin{proofbox}
\textbf{方法一：由第一定义作和差，避免根式展开（优先）。}
设 $P(x,y)$ 在椭圆上，记 $r_L=PF_L,r_R=PF_R$。
第一定义给出 $r_L+r_R=2a$，距离平方差给出 $r_L^2-r_R^2=4cx$。
两式相除后解得
\[
r_R=a-\frac ca x=a-ex,\qquad r_L=a+ex.
\]
又因为 $a^2/c>a\ge x$，右准距为 $a^2/c-x$，所以
\[
e\,d(P,\ell_R)=\frac ca\left(\frac{a^2}{c}-x\right)=a-ex=r_R.
\]
左边同理。这里准线位置正是把焦半径 $a-ex$ 写成 $e$ 乘一个水平距离所得。

\textbf{方法二：平方差恒等式，一次核验两个方向。}
对任意平面点 $P(x,y)$，利用 $c^2=a^2-b^2$ 直接展开：
\[
\begin{aligned}
PF_R^2-e^2d(P,\ell_R)^2
&=(x-c)^2+y^2-\frac{c^2}{a^2}(x-a^2/c)^2\\
&=b^2\left(\frac{x^2}{a^2}+\frac{y^2}{b^2}-1\right).
\end{aligned}
\]
右端为零恰是椭圆方程；左端为零又等价于 $PF_R=e\,d(P,\ell_R)$，因为两距离和 $e$ 都非负。
若准距为零，则此等式会迫使 $P=F_R\in\ell_R$，矛盾，故能除以准距。
这就补齐了反向验证，而非只证椭圆上的点满足比例。
\end{proofbox}
\clearpage
\begin{proofbox}
\textbf{一般焦点与准线如何得到椭圆？}
取焦点在准线上的垂足为原点，准线为 $x=0$，焦点为 $(q,0)$，$q>0$。
条件等价于
\[
(x-q)^2+y^2=e^2x^2.
\]
因 $0<e<1$，配方得
\[
(1-e^2)\left(x-\frac{q}{1-e^2}\right)^2+y^2
=\frac{q^2e^2}{1-e^2}.
\]
右端为正，两个平方项系数均正，除以右端即得到平移后的椭圆标准方程。
半轴为 $a=qe/(1-e^2)$，$b=qe/\sqrt{1-e^2}$，半焦距 $c=qe^2/(1-e^2)$，故 $c/a=e$。
由于 $e>0$，原准线 $x=0$ 上没有点满足距离等式，距离比始终有意义。
\end{proofbox}
